\section{Related work}\label{sec:related-work}
\textcolor{red}{TODO: describe diffusion models as in between?}
\textcolor{red}{
    TODO: IMPORTANT! describe everything in a view that the model solves isomorphism OUTSIDE the model and not inside.
    make it a solution rather than a problem shift since all other methods solve iso in the model. write it in a way
    that sees the transfer as analogy between one-shot and sequential generators, ie, introducing a sequential algorithm class for
    exact decoding similar as for generative models.
}
\textcolor{red}{
    TODO: section should review not autoencoder but generators of graphs, where decoders are one special case.
}\\
\textcolor{red}{TODO: even GRALE is not 100\% unsupervised since labels are generated for matching}\\
\textcolor{red}{TODO: Molecule generators: this~\cite{mercadopractical}.}\\
\textcolor{red}{TODO: state that sequential generation is superior since isomorphism is a harder problem to solve than iterative decoding}\\

In this paper, we study the decoding of graphs from vector representations, e.g., a graph-level embedding produced
by an encoder or obtained from a fixed algorithm, to the actual graph.
This setting is distinct from \emph{unconditional} graph generative modeling, in which the objective is to learn to generate
graphs similar to those of a reference distribution.
However, as both concepts need to derive procedures for generating the graph, it serves as the conceptual foundation of decoder
design, in particular regarding one-shot versus sequential construction and mechanisms for handling graph representation and isomorphism
as well as validity constraints, if any.

\paragraph{One-shot vs. sequential generation.}
Generative models for graphs primarily aim to model an underlying distribution with the ultimate goal
of being able to sample new items following this distribution.
They focus less on representing single samples exactly.
One-shot methods such as variational graph autoencoders~\cite{vgae,graphvae} represent the output graph as a single object
in Euclidean space and optimize the graph with point-wise losses for reconstruction, similar to images.
While these approaches are computationally efficient and simple to implement with feedforward architectures,
they inherently output an explicitly ordered representation of the graph and thus need to account for graph isomorphism during training.
Since graph isomorphism is a challenging problem to solve, another line of research focuses on generating graphs
sequentially~\cite{graphrnn,gran,learning-deep-gen-graphs,grover2019graphite}.
Generating the graph one element at a time, these models learn to predict the next step recursively by taking into account
partial results and as such, make the node order implicit and do not suffer from graph isomorphism.
However, sequential methods must devise a generation procedure involving several forward passes for inference as well as step-wise supervision
of the process for training.
Their recursive nature makes generation inherently slower than one-shot generation but allows for more flexible and controlled generation.


\paragraph{Unconditional vs. conditional generation.}
Besides algorithmic differences in how to generate a graph, there are also differences in the objective of graph generation,
As generative modeling aims to generate graphs representative of a reference distribution and without restrictions of
which graph to generate, we call this branch \emph{unconditional generation}.
This form of generation is mature and well-studied, and the earliest approaches date back to ER and BA graphs\textcolor{red}{(cite)}.
\textcolor{red}{Explain modern neural network versions}.
As unconditional graph generation is rather unrestricted, these methods scale to large unattributed graphs~\cite{graphrnn,gran}
or smaller attributed graphs, e.g., small molecules~\cite{learning-deep-gen-graphs}.
However, these methods are unable to generate a specific graph given their unconditional nature.
In contrast, powerful methods for conditional graph generation may be used in a wide spectrum of tasks, e.g.,
whenever the underlying task is to translate an input to a graph.
For example, a model generating graphs conditioned on an image embedding can solve a wide array of problems, including
the extraction of graphs from an image.
As such, another line of research focuses on building conditional graph generators.

Building on the seminal work of variational graph autoencoders~\cite{vgae}, which only reconstruct the graph
based on their given nodes, the authors of~\cite{graphvae} propose the first graph-level extension.
They recognize that graph isomorphism is the inherent problem of reconstructing graphs in the autoencoder framework
and propose to introduce graph matching to find correspondences and to calculate losses pointwise as in regular
autoencoders.
In~\cite{winter2021pigvae}, the authors introduce a permutation-invariant graph variational autoencoder (PIGVAE).
This work goes one step further and establishes an architectural design in which the model learns a canonical ordering internally.
The authors explain it with the argument that during global aggregation, the node order is lost in the graph latent.
As such, a model needs to learn to infer it, which they solve by introducing an additional predictor of graph size from the latent
as well as a permuter model that predicts the alignment.
More recently, a line of research has been proposed under the term of Supervised Graph Prediction (SGP), in which the authors
go one step further and develop some of the first methods that directly translate different modalities into
graphs~\cite{krzakala2024framework,any2graph2024, grale2025}.
In~\cite{any2graph2024}, the authors seek to develop a permutation-invariant distance for graphs and introduce a loss based
on Optimal Transport that is permutation invariant and handles variable-sized inputs.
A notable improvement is that their model, Any2Graph, can decode arbitrary embeddings of different input modalities into their corresponding
graph representation.
In follow-up work~\cite{grale2025}, the authors combine the loss introduced for Any2Graph and the matching strategy of PIGVAE.
The resulting architecture, GRALE, operates on both graph and node embeddings, where the graph-level embeddings serve as a latent
representation and node-level embeddings are used for matching, achieving state-of-the-art results.

%
%In summary, graph autoencoders aim to propose loss functions that \ldots
%
%
%\begin{enumerate}
%    \item \cite{jtvae} (\citeyear{jtvae})
%    \begin{itemize}
%        \item contracts graph into an element tree, need pre-processing
%        \item unclear why this needs to be variational
%        \item learning predicts the connections of subgraphs to the tree
%        \item \textbf{encoding and decoding is stochastic, our method is exact}
%        \item \textbf{scoring remotely similar, but subgraph decoding does not have all this model machinery}
%        \item \textbf{\red{they use ZINC and achieve a reconstruction accuracy of 76.7\%, this is easily beatable}}
%        \item also the argument that atom-by-atom generates invalid results is not always true
%    \end{itemize}
%
%    \item \cite{jin2020hierarchical} (\citeyear{jin2020hierarchical})
%    \begin{itemize}
%        \item uses larger motifs, depends on extracting them, paper extracts 436 different ones
%    \end{itemize}
%
%    \item \cite{kearnes2019decoding} (\citeyear{kearnes2019decoding})
%    \begin{itemize}
%        \item uses RL, achieved only 66\% on QM9, also reports that \cite{jtvae} also does not achieve much higher accuracies
%    \end{itemize}
%
%
%\end{enumerate}
